\documentclass{article}
\usepackage[usenames]{color}
\usepackage{palatino}
\usepackage[all]{xy}
\usepackage[toc,page]{appendix}
\usepackage{tocloft}
\usepackage[french]{babel}
\usepackage[T1]{fontenc}
\usepackage[utf8]{inputenc}

\begin{document}
\begin{center}\textsc{Citations}\end{center}
\vspace{30pt}

\begin{center} « C’est en vain qu’on court après l’existence et la vérité. Tout est néant, une ronde d’hallucinations sans rythme. Ce qui fait qu’une chose est, c’est notre état de fièvre, et la vivacité de nos ardeurs projette des vérités sur un monde d’absences. Le souffle de la substance qui transforme en réalité le non-être du monde émane de nos intensités. Si nous étions plus froids et plus calmes, rien ne serait. Les deux intérieurs soutiennent la solidité apparente du monde, animent le décor vide et sont les véritables architectes de la vie. Le monde est une prolongation extérieur de nos flammes. » \end{center} 

\begin{flushright}\textit{Le crépuscule des pensées}, Emile Cioran\end{flushright}
\vspace{8pt}
\begin{center} « Car jusqu’à quand laisserons nous piétiné par tant de créateurs qui n’ont rien su, enfants terribles et inspirés, dépourvus de la maturité du bonheur et du malheur? Un génie qui n’a pas atteint les racines de la vie, quelle que soit ses multiples possibilités d’expression, ne doit être goûter qu’aux instants d’indifférence. Il est plus horrible de penser que ceux qui ont vraiment su quelque chose soient si rares, que le nombre des existences complètes soit encore restreint. Mais qu’est ce qu’une existence complète et qu’est ce que savoir ? — garder une soif de la vie dans les crépuscules... »\end{center}

\begin{flushright}\textit{Le crépuscule des pensées}, Emile Cioran\end{flushright}
\vspace{8pt}
\begin{center} « Il regarda les femmes et les servantes se rendant au marché aux poissons, les yeux fixés sur poissonniers et leurs épouses mal embouchées qui offraient et vantaient leur marchandise, sortaient la bouche douloureusement ouverte, leur yeux d’or anxieux et fixes, s’abandonnaient avec résignation à la mort ou de débattaient contre elle avec fureur et désespoir. Comme il était arrivé maintes fois déjà il fut pris de pitié pour ces animaux et de dégoût à l’égard des hommes; pourquoi étaient-ils si insensibles et brutaux, si immensément bêtes et stupides; pourquoi ne voyaient-ils rien, tous, les poissonniers et les poissonnières, et les clients qui marchandaient, pourquoi ne voyaient-ils pas ces bouches, ces yeux dont l’angoisse de mort, ces queues frappant furieusement autour d’elles, cette affreuse et inutile lutte désespérée, cette transformation intolérable des animaux mystérieux, merveilleusement beaux, dernier léger frisson de la mort passant sur leur peau agonisante avant qu’ils soient la, allongés, morts et éteints, lamentables morceaux de viande destinés à la table des mangeurs réjouis? » \end{center}

\begin{flushright}\textit{Narcisse et Goldmund}, Hermann Hesse\end{flushright}
\vspace{7pt}
\begin{center} « Dieu mort, restent les hommes, c’est-à-dire l’histoire qu’il faut comprendre et bâtir. Le nihilisme, qui, au sein de la révolte, submerge alors la force de la création, ajoute seulement qu’on peut la bâtir par tous les moyens. Aux cimes de l’irrationnel, l’homme, sur une terre qu’il sait désormais solitaire, va joindre les crimes de la raison en marche vers l’empire des hommes. Au « je me révolte, donc nous sommes », il ajoute, méditant de prodigieux desseins et la mort même de la révolte : « et nous sommes seuls. » » \end{center}

\begin{flushright}\textit{L'Homme révolté}, Albert Camus\end{flushright}
\vspace{7pt}
\begin{center} « Les sociétés modernes souffrent d’un manque cruel de politique, comblé par la gestion au jour le jour des revendications particulières, la police de l’espace publique et la mise en musique des impératifs comptables au services de projets souvent improvisés ou mal pensés. C’est aussi l’une des raisons du constat répandu que les « mots » n’ont plus de sens ; le recours aux éléments de langage constitue la forme achevée d’un mépris consommé pour la vérité et l’expression d’une institutionnalisation cynique du mensonge. » \end{center}

\begin{flushright}\textit{Disparition de l’espace du débat}, Le monde diplomatique\end{flushright}
\vspace{7pt}
\begin{center} « Ich werde mich dann zur Mathematik wenden, und Auskunft über Sinn und Ursprung der Symbole zu erhalten, und wir werden da den Menschen, sofern er schöpferischer Geist ist, als den Baumeister der Symbolwelt entdecken. Der Geist ist Freiheit in der Gebundenheit seines Daseins; sein Feld ist das Mögliche, das geöffnet ist zum Unendlichen hin, im Gegensatz zum geschlossen Sein.Nur indem die Freiheit des Geistes sich selber bindet an das Gesetz, begreift der Geist nachkonstruirend die Gegebenheit der Welt und seines eigenen Daseins in der Welt. » \end{center}

\begin{flushright}Hermann Weyl, 1948\end{flushright}
\vspace{7pt}
\begin{center} « Suis ton plan, cher Lucilius ; reprends possession de toi-même : le temps qui jusqu’ici t’étais ravi, ou dérobé, ou que tu laissais perdre, recueille et ménage-le. Persuade-toi que la chose a lieu comme je te l’écris : il est des heures qu’on nous enlève par force, d’autres par supplice, d’autres coulent de nos mains. Or la plus honteuse perte est celle qui vient de négligence ; et, si tu y prends garde, la plus grande part de la vie se passe à mal faire, une grande à ne rien faire, le tout à faire autre chose que ce qu’on devrait. Montre-moi un homme qui mette au temps le moindre prix, qui sache ce que vaut un jour, qui comprenne que chaque jour il meurt en détail ! Car c’est notre erreur de ne voir la mort que devant nous : en grande partie déjà on l’a laissée derrière : tout l’espace franchi est à elle. » \end{center}

\begin{flushright}\textit{Lettre I, Sur l’emploi du temps}, in \textit{Lettres à Lucilius}, Sénèque\end{flushright}
\vspace{8pt}
\begin{center} « Dans ce milieu à température constante, dans ce climat préservé des catastrophes, des orages, la Raison a poussé comme une plante solitaire, identique sous une diversité apparente qui ne trompa jamais les initiés. [...] Le monde qui est l’objet de la philosophie est une construction des techniques, des sciences et des actions. [...] le principe de l’explication des philosophes repose dans les changements du monde humain et des milieux de l’intelligence. Les faits de l’homme commandent les relations de la pensée avec ses objets, le développement des groupes humains commande le développement de toutes les activités techniques, politiques, sociales est le moteur de ce que les penseurs nomment le mouvement de l’Esprit. » \end{center} 

\begin{flushright}\textit{Les chiens de garde}, Paul Nizan\end{flushright}
\vspace{6pt}
\begin{center} « Le grand thème du Temps retrouvé est celui-ci : la recherche de la vérité est l’aventure propre de l’involontaire. La pensée n’est rien sans quelque chose qui force à penser, qui fait violence à la pensée. [...] ce qui force à penser, c’est le signe. Le signe est l’objet d’une rencontre; mais c’est précisément la contingence de la rencontre qui garanti la nécessité de ce qu’elle donne à penser. L’acte de penser ne découle pas d’une simple possibilité naturelle. Il est au contraire la seule création véritable. Nous ne cherchons la vérité que dans le temps, contraints et forcés. » \end{center} 

\begin{flushright}\textit{Proust et les signes}, Gilles Deleuze\end{flushright}
\vspace{6pt}
\begin{center} « Mais « politique », « économie », ce ne sont là que des maux dérivés. L’essentiel est ailleurs. il est d’abord dans la manie fondationnelle - constituante. Car fonder, c’est proférer le mensonge inaugural : « le mensonge de la civilisation consiste à couvrir de silence ses multiples décisions, puis à s’inventer une histoire logique. ». Les groupes ne sortent du chaos de la condition anarchique que par un geste d’ancrage arbitraire, l’ancrage n’étant soutenable qu’à la condition que l’ordre ainsi institue institue dans le même temps la méconnaissance de son propre arbitraire. » \end{center}

\begin{flushright}\textit{Vivre sans ?}, Frédéric Lordon\end{flushright}
\vspace{6pt}
\begin{center} « Et le plus grave, la tentation à laquelle tous succombent volontiers, Nietzsche, mais plus sournoisement Pascal ou Wittgenstein : pouvoir mépriser les mathématiques, les réduire, au regard du sérieux moral, de l’intensité existentielle, à un jeu d’enfant. Le mépris des mathématiques est ce dont aucune philosophie ne se relève. C’est pourquoi j’arrête l’oeuvre de Wittgenstein au Tractatus, ou ce mépris est déjà là, mais mêlée, encore, d’adoration. » \end{center} 
\begin{flushright}\textit{L'antiphilosophie de Wittgenstein}, Alain Badiou\end{flushright}
\vspace{7pt}
\begin{center} « Un Wittgenstein philosophique, qui s’efforce, sans jamais y parvenir réellement, d’atteindre à un idéal moral de type spinoziste. […] Le refus de moraliser la morale, l’éternité conçue comme la vie dans le présent, le refus de juger et se juger, l’acceptation de ce sui arrive et de ce qui nous arrive comme la seule possibilité réelle d’affranchissement à l’égard du destin, la recherche de béatitude personnel dans la solitude, la pauvreté et l’ascèse. » \end{center} 

\begin{flushright}\textit{Wittgenstein, La rime et la raison}, Jacques Bouveresse\end{flushright}
\vspace{7pt}
\begin{center} « « Les révoltes logiques », la philosophie est la combinaison d’un désir de révolution – le bonheur réel exigeant qu’on se lève contre le monde tel qu’il est – et d’une exigence de la rationalité – la pulsion révoltée ne pouvant à elle seule parvenir aux buts qu’elle se fixe.
Le désir philosophique est bien, de façon extrêmement générale, le désir d’une révolution dans la pensée, aussi bien collective que personnelle, et ce en vue d’un bonheur réel distingué du semblant de bonheur qu’est la satisfaction.

La maxime pourrait être : en finir avec la fin. Et en finir suppose qu’une décision soit prise. » \end{center} 

\begin{flushright}\textit{Métaphysique du bonheur réel}, Alain Badiou\end{flushright}
\vspace{7pt}
\begin{center} « Les structures et les institutions n’ignorent pas les psychés, pour cette simple raison qu’elles opèrent à travers elles — comment, sans quelque relais dans les individus, pourraient-elles agir sur les individus, et leur faire faire quoi que ce soit ? Ou alors il faudra considérer que les « nécessités de l’ordre capitaliste », par exemple, s’accomplissent toutes seules, que les individus sont mus à leur service par une opération miraculeuse, qui restera un mystère. En réalité, les institutions œuvrent par les psychés. Mais elles font davantage encore : elles sélectionnent les psychés, disons plus précisément : les complexions psychiques, les plus adéquates à leurs propres réquisits. L’Église sélectionne ses psychés adéquates. La police sélectionne ses psychés adéquates — adéquates à ses finalités d’institution de la violence. L’entreprise néolibérale sélectionne ses psychés adéquates – des pervers, au sens clinique du déni de l’altérité, condition de l’instrumentation des humains comme des choses. » \end{center} 

\begin{flushright}\textit{Sont-ils fous ?}, Frédéric Lordon\end{flushright}
\vspace{8pt}

\begin{center} « Il est crucial, à mes yeux, pour un enfant, d’être exposé très tôt à la musique. Je pense qu’exposer un enfant à la musique, vers l’âge de cinq ou six ans, permet d’équilibrer un petit peu la prépondérance dans son intellect du sens de la vue, de cette richesse incroyable purement visuelle, qu’un enfant acquiert très tôt et qui donc, en fait, est reliée à la géométrie. La musique permet de l’équilibrer par l’algèbre, c’est-à-dire qu’elle s’inscrit dans le temps, exactement comme l’algèbre s’inscrit dans le temps. Dans les mathématiques il y a cette dualité fondamentale entre d’un côté la géométrie, qui correspond aux aires visuelles du cerveau, et qui donne une intuition instantanée, immédiate. On voit une figure géométrique, boum! C’est ça, c’est tout, on n’a même pas besoin d’expliquer, on n’a pas envie d’expliquer. Et d’un autre côté, il y a l’algèbre. L’algèbre, cela n’a rien de visuel, en revanche, cela a une temporalité, ça s’inscrit dans le temps ! C’est le calcul, etc. C’est quelque chose qui évolue, et c’est quelque chose qui est très proche du langage et qui donc a la précision diabolique du langage. Et l’on peut percevoir cette puissance, l’élaboration de l’algèbre, à travers la musique. Donc, pour moi, il y a une connivence incroyable, justement, entre la musique perçue comme cela, et l’algèbre. Par exemple, j’adore certains préludes de Chopin parce que je trouve qu’ils ont exactement cette merveilleuse propriété de condensation, de distillation. C’est une musique qui arrive dans une pièce un petit peu comme si la fenêtre s’ouvrait brutalement par un coup de vent, et puis repart de l’autre côté. Condenser une idée sous sa forme la plus limpide, sous la forme la plus pure qui soit... L’algèbre, c’est ça, d’une certaine manière. » \end{center} 

\begin{flushright}\textit{Algèbre et musique}, Alain Connes\end{flushright}
\vspace{8pt}
\begin{center} « Toute réalité humaine est une passion en ce qu’elle projette de se perdre pour fonder l’être et pour constituer de même coup l’en-soi qui échappe à la contingence en étant son propre fondement, ens causa sui que les religions nomment Dieu. Ainsi, la passion de l’homme est-elle inverse de celle du Christ, car l’homme se perd en tant qu’homme pour que Dieu naisse. Mais l’idée de Dieu est une contradiction et nous nous perdons en vain. L’homme est une passion inutile.» \end{center} 

\begin{flushright}\textit{L’Être et le Néant}, Jean-Paul Sartre\end{flushright}
\vspace{8pt}
\begin{center} « Tout en niant l’existence du vide et en soutenant l’identification cartésienne entre étendue et matière, Spinoza distingue soigneusement l’étendue, telle qu’elle est donnée aux sens et représentée par l’imagination, et celle perçue par l’entendement. La première, divisible et susceptible de mouvement (correspondant au monde cartésien indéfiniment étendu) constitue l’éternelle multiplicité des modes perpétuellement changeants et finis, la seconde, véritablement et pleinement infinie et par conséquent indivisible, constitue l’attribut essentiel et éternel de l’Être existant a se et per se, c’est-à-dire Dieu. » \end{center} 

\begin{flushright}\textit{Du monde clos à l’univers infini}, Alexandre Koyré\end{flushright}
\vspace{8pt}
\begin{center} « Cavaillès note la faiblesse philosophique d’une épistémologie qui croit pouvoir caractériser les sciences comme des systèmes hypothético-déductifs : « Comment un principe ou une réunion de principes qui, dans leur contenu et dans leur rassemblement, ne sont pas eux même intelligibles, peuvent-ils être point de départ pour un déroulement intelligible ? L’alliance hétérogène d’un concret pur constaté et d’un enchaînement rationnel est simple image sans pensée. » 
	Le problème est donc : d’appréhender le processus créateur, en retrouver la « structure , non par description, mais apodictiquement en tant qu’elle se déroule et se démontre elle-même. Autrement dit, la théorie de la science est un a priori, non antérieur à la science, mais âme de la science. » » \end{center} 
	
\begin{flushright}\textit{Jean Cavaillès, un philosophe dans la guerre}, Gabrielle Ferrières\end{flushright}
\vspace{8pt}
\begin{center} « On oppose souvent deux formes de pensée critiques. L’une, celle de Marx et Horkheimer, considère que nous sommes prisonniers d’une idéologie qui nous empêche de voir le monde tel qu’il est : le travail critique consiste donc à nous émanciper par le dévoilement de cette aliénation, c’est-à-dire des logiques, des intérêts et des pouvoirs qui font obstacle à l’analyse de notre condition — encore que les dominés fassent plus souvent preuve l’intelligence sociale que la théorie de ne suppose. L’autre, celle de Nietzsche et de Wittgenstein, affirme que nous sommes captifs d’une vision des choses qui nous les faits voir faussement apparaître comme évidentes : le travail critique consiste alors à nous rendre conscients du caractère arbitraire et contingent de nos valeurs, de nos normes et de nos représentations, autrement dit à remettre en cause ce que la force et l’habitude et l’empire du conformisme nous font tenir pour acquis — mais là aussi il n’est d’usage de le croire. » \end{center} 

\begin{flushright}\textit{La vie, mode d’emploi critique}, Didier Fassin\end{flushright}
\vspace{8pt}
\begin{center} « Dans une lettre à R. Hädrich du 8 septembre 1954, à propos d’une adaptation de L’Etranger pour la scène allemande, Camus fait la critique d’un « monde » sous la domination du « jugement » : « Ce qui est attaqué de front, ici, ce n’est pas la morale mais le monde du procès, qui est aussi bien bourgeois que nazi ou communiste, qui en un mot le chancre contemporain […]. Meursault n’est pas du côté des juges, de la loi sociale, des sentiments convenus. Il existe, comme une pierre, ou le vent, ou la mer sous le soleil, qui, eux, ne mentent jamais. Si vous envisagez le livre sous cet aspect, vous y verrez une morale de la sincérité et une exaltation à la fois ironique et tragique à la joie du monde. Ce qui exclut l’ombre, la caricature expressionniste ou la lumière désespérée ». 
	Le monde du jugement et le monde de la production sont un seul et même monde dont l’« unité est celle du nihilisme ». Et Meursault est la puissance active de leur subversion.
» \end{center}

\begin{flushright}\textit{Albert Camus, de la transfiguration}, Laurent Bove\end{flushright}
\vspace{8pt}
\begin{center} « Hum, oui, toutes les choses sont à la portée de l'homme, et tout lui passe sous le nez, à cause de sa poltronnerie... c'est devenu un axiome... Il serait curieux de savoir ce que les hommes redoutent pardessus tout. Ce qui les tire de leurs habitudes, voilà ce qui les effraye le plus... Mais je bavarde beaucoup trop c'est pourquoi je ne fais rien, ou peut-être devrais-je dire que c'est parce que je ne fais rien que je bavarde. Ce mois-ci j'ai pris l'habitude de monologuer, couché pendant des jours entiers dans mon coin, à songer... à des sottises. [...] Suis-je vraiment capable de "cela" ? Un jeu ! oui, c'est bien cela, un jeu ! » \end{center}

\begin{flushright}Raskolnikov dans \textit{Crimes et Châtiments}, Fiodor Dostoïevski\end{flushright}
\vspace{8pt}
\begin{center} « Il est encore heureux que vous n'avez tué cette vieille, mais, avec une autre théorie, vous auriez pu commettre une action cent millions de fois pire. Remerciez donc Dieu de ne pas l'avoir permis, car Il a peut-être — qui le sait ? — des desseins sur vous. Et vous, vous avez du courage, ne reculez pas, par pusillanimité, devant la grande action qu'il vous reste à accomplir. Il serait honteux pour vous d'être lâche. Si vous avez commis l'acte, eh bien, soyez fort et faites ce qu'exige la justice. Je sais que vous ne me croyez pas, mais je vous donne ma parole que vous reprendrez goût à la vie. En ce moment, il ne vous faut que de l'air, de l'air, de l'air... » \end{center}

\begin{flushright}Porphyre dans \textit{Crimes et Châtiments}, Fiodor Dostoïevski\end{flushright}
\vspace{8pt}
\begin{center} « J'ai compris, Sonia, qu'à vouloir attendre que tout le monde devienne intelligent, on risque de perdre beaucoup de temps... Ensuite, j'ai pu me convaincre que ce moment n'arriverait jamais, que les hommes ne pouvaient changer. L'essayer n'eût été qu'une perte de temps inutile. C'est la loi humaine... la voie, Sonia, voilà ! Et maintenant, je sais que celui qui est doué d'une volonté, d'un esprit puissants, n'a pas beaucoup de peine à devenir leur maître. Qui ose beaucoup a raison devant eux. Qui les brave et méprise gagne leur respect. Il devient leur législateur. » \end{center}

\begin{flushright}Raskolnikov dans \textit{Crimes et Châtiments}, Fiodor Dostoïevski\end{flushright}
\vspace{8pt}
\begin{center} « La liberté sera totale quand il sera indifférent de vivre ou de mourir. Voilà le but de tout. [...] La vie est souffrance, la vie est terreur, et l'homme est malheureux. Maintenant l'homme aime la vie parce qu'il aime la souffrance et la terreur. Aujourd'hui l'homme n'est pas encore l'homme. Un homme nouveau viendra, heureux et fier. Celui auquel il sera indifférent de vivre ou de ne pas vivre, celui-là sera l'homme nouveau. Celui qui vaincra la souffrance et la terreur, celui-là sera lui-même Dieu. Quant à l'autre Dieu, il ne sera plus. » \end{center}

\begin{flushright}Kirilov dans \textit{Les Démons}, Fiodor Dostoïevski\end{flushright}
\vspace{8pt}
\begin{center} « La vérité qui, selon Jésus, nous rendra libre n'est pas la vérité objective des sciences, et elle n'est pas non plus la vérité de la théologie : pas plus qu'elle n'est un ouvrage de cosmologie, la Bible n'est pas un manuel d'anthropologie ou de théologie. La seule vérité que nous révèle l'Ecriture, celle qui au cours du temps, ne peut subir aucune démystification – puisqu'elle n'est pas un énoncé expérimental, logique, métaphysique, mais un appel pratique – c'est la vérité de l'amour, de la \textit{caritas}. » \end{center}

\begin{flushright}cité dans \textit{Peut-on ne pas croire ?}, Jacques Bouveresse\end{flushright}
\begin{center} « La vérité sort plutôt de l'erreur que de la confusion. » \end{center}

\begin{flushright}Ernest Renan cité dans \textit{Peut-on ne pas croire ?}, Jacques Bouveresse\end{flushright}
\vspace{8pt}
\begin{center} « L'amour, c'est en vérité la distance intérieure de l'individu au lien de société, sa déliaison intime de l'originelle obligation de communauté. Non pas leur rejet ou leur contestation ouverte : leur désertion silencieuse et privée. Jésus ne déclare pas dissous l'ancien principe d'appartenance et le système des devoirs qui l'accompagne. Il désigne et instaure à côté une tout autre compréhension de l'obligation, à base d'autonomie du cœur. Il ne fonde pas un ordre individualiste. Mais il suscite une personne, un individu du dedans entièrement dégagés, du point de vue de la règle à laquelle ils obéissent, et des fins qu'ils poursuivent, de la loi de l'inclusion qui régit ce bas monde.
Inversion de \textit{doctrine}, donc ; inversion, enfin, de \textit{destin}. Jésus ne va pas au triomphe, il va à la \textit{défaite}. »\end{center}

\begin{flushright}\textit{Le désenchantement du monde}, Marcel Gauchet\end{flushright}
\vspace{8pt}
\begin{center} « Pas plus que l'immanence de la création, le travail dans cette immanence n'est imaginable. Les possibilités de la puissance sont inconnues. L'homme même ne connaît pas la puissance de l'home. Le travail humain est déjà une telle force transfigurante qu'on ne peut songer sans éblouissement au travail divin. une femme pelure, chimiste Smithson est là, il recueille une goutte de cette eau, et cette larme d'une femme devient une formule chimique, d'où sortira une branche de la science. Byron achète un schilling chez son papetier une bouteille d'encre qu'il revendra cent mille francs à son éditeur. Et nous nous bornons à l'énoncer du résultat matériel ; le résultat oral est bien plus surprenant encore. Un certain travail, s'infiltrant à une masse de métal ou de pierre, à une toile, à une feuille de papier, lui fait subir une telle sublimation que de matière qu'elle était, elle devient une idée. Du travail sort une métaphysique, réfractaire à toute formule, productrice de forces et de valeurs. La mise en oeuvre est une seconde création. La première création n'est que la mise en mouvement. Après la sève, l'intelligence. [...] Si les Prométhées d'en bas, surprenant et dérobant au créateur son secret, peuvent de telles choses, s'ils les réalisent, que ne réaliseront pas les providences d'en haut ! que ne pourra point ne créateur lui-même ? \textit{Quid domini facient, audent cum talia fures !} »\end{center}

\begin{center} « Il y a dans la création un Inconnu. Cet Inconnu a ses raisons. Son pourquoi nous déborde. Il se dépense dans l'effroi comme dans la splendeur. Ses réussites dans le terrible font frémir. Le rêve de l'homme est un essai toujours dépassé par la création ; il y a quelque chose de plus nocturne que le songe, c'est le fait. La réalité distance le cauchemar. [...] Mais l'imagination renonce à plonger et à planer; la science refuse de tâtonner. Une connaissance trop profonde de la réalité dans un vivant paraîtrait aux autres vivants folie, la science complète aurait un air de démence, et l'infortuné qui parviendrait à se rencontrer face à face avec le Grand Inconnu sur le sommet des choses ne redescendrait du Sinaï que pour entrer à Bedlam. » \end{center}

\begin{flushright}Reliquats des \textit{Travailleurs de la mer}, Victor Hugo\end{flushright}
\vspace{8pt}

\begin{center} « Dans la morale, comme dans l'art, dire n'est rien, faire est tout. L'idée qui se cache sous un tableau de Raphaël est peu de chose; c'est le tableau seul qui compte. De même, en morale, la vérité ne prend quelque valeur que si elle passe à l'état de sentiment, et elle n'atteint tout son prix que quand elle se réalise dans le monde à l'état de fait. Des hommes d'une médiocre moralité ont écrit de fort bonnes maximes. Des hommes très-vertueux, d'un autre côté, n'ont rien fait pour continuer dans le monde la tradition de la vertu. La palme est à celui qui a été puissant en paroles et en œuvres, qui a senti le bien, et au prix de son sang l'a fait triompher. Jésus, à ce double point de vue, est sans égal; sa gloire reste entière et sera toujours renouvelée. » \end{center}

\begin{flushright}\textit{Vie de Jésus}, Ernest Renan\end{flushright}

\begin{center}« Jésus ne sera plus seulement un délicieux moraliste, aspirant à renfermer en quelques aphorismes vifs et courts des leçons sublimes; c'est le révolutionnaire transcendant, qui essaye de renouveler le monde par ses bases mêmes et de fonder sur terre l'idéal qu'il a conçu. « Attendre le royaume de Dieu » sera synonyme d'être disciple de Jésus 2. Ce mot de « royaume de Dieu » ou de « royaume du ciel, » ainsi que nous l'avons déjà dit 3, était depuis longtemps familier aux Juifs. Mais Jésus lui donnait un sens moral, une portée sociale que l'auteur même du Livre de Daniel, dans son enthousiasme apocalyptique, n'avait pas même osé entrevoir.

Dans le monde tel qu'il est, c'est le mal qui règne. Satan est le « roi de ce monde, » et tout lui obéit. Les rois tuent les prophètes. Les prêtres et les docteurs ne font pas ce qu'ils ordonnent aux autres de faire. Les justes sont persécutés, et l'unique partage des bons est de pleurer. Le «monde» est de la sorte l'ennemi de Dieu et de ses saints 1; mais Dieu se réveillera et vengera ses saints. Le jour est proche; car l'abomination est à son comble. Le règne du bien aura son tour. »
\end{center}

\begin{center}« Une révolution radicale, embrassant jusqu'à la nature elle-même, telle fut donc la pensée fondamentale de Jésus. Dès lors, sans doute, il avait renoncé à la politique; l'exemple de Juda le Gaulonite lui avait montré l'inutilité des séditions populaires. Jamais il ne songea à se révolter contre les Romains et les té-traques. Le principe effréné et anarchique du Gaulonite n'était pas le sien. Sa soumission aux pouvoirs établis, dérisoire au fond, était complète dans la forme. Il payait le tribut à César pour ne pas scandaliser. La liberté et le droit ne sont pas de ce monde; pourquoi troubler sa vie par de vaines susceptibilités?
Méprisant la terre, convaincu que le monde présent ne mérite pas qu'on s'en soucie, il se réfugiait dans son royaume idéal; il fondait cette grande doctrine du dédain transcendant, vraie doctrine de la liberté des âmes, qui seule donne la paix. Mais il n'avait pas dit encore: « Mon royaume n'est pas de ce monde. »

Bien des ténèbres se mêlaient à ses vues les plus droites. Parfois des tentations étranges traversaient son esprit. Dans le désert de Judée, Satan lui avait proposé les royaumes de la terre. »\end{center}

\begin{center}« Ce qui distingue, en effet, Jésus des agitateurs de son temps et de ceux de tous les siècles, c'est son parfait idéalisme. Jésus, à quelques égards, est un anarchiste, car il n'a aucune idée du gouvernement civil. Ce gouvernement lui semble purement et simplement un abus. Il en parle en termes vagues et à la façon d'une personne du peuple qui n'a aucune idée de politique. »\end{center}

\begin{center}« L'idée qu'on est tout-puissant par la souffrance et la résignation, qu'on triomphe de la force par la pureté du cœur, est bien une idée propre de Jésus. Une immense révolution sociale, où les rangs seront intervertis, où tout ce qui est officiel en ce monde sera humilié, voilà son rêve. Le monde ne le croira pas; le monde le tuera. Mais ses disciples ne seront pas du monde. Ils seront un petit troupeau d'humbles et de simples, qui vaincra par son humilité même.
Le sentiment qui a fait de « mondain» l'antithèse de « chrétien » a, dans les pensées du maître, sa pleine justification. »\end{center}

\begin{center}« Mais moi je sais, (dit Dieu), que rien n’est grand comme l’ordre et que rien n’est profond comme le labour. »\end{center}

\begin{flushright}\textit{Le porche du mystère de la deuxième vertu}, Charles Péguy\end{flushright}

\begin{center}« Et certainement la plupart des erreurs consistent seulement en ceci que nous n'appliquons pas correctement les noms aux choses. »\end{center}

\begin{flushright}\textit{L'Ethique}, 2.47 scolie, Spinoza\end{flushright}


\begin{center}« Du troisième genre de connaissance naît nécessairement un Amour intellectuel de Dieu. Car de ce genre de connaissance naît (par la Prop. précédente) une Joie qu’accompagne l’idée de Dieu comme cause, […] »\end{center}

\begin{flushright}\textit{L'Ethique}, 5.32 corollaire, Spinoza\end{flushright}

\begin{center}« Et certes les passions humaines marquent l’art et la puissance de la nature, sinon celle de l’homme, non moins que beaucoup d’autres choses que nous admirons et dont la contemplation nous enchante. » \end{center}

\begin{flushright}\textit{L'Ethique}, 4.47 scolie, Spinoza\end{flushright}


\begin{center}« Comme vous connaissez mal le bonheur des hommes, gens du confort et de bonne volonté ! – car le bonheur et le malheur sont des frères jumeaux qui grandissent ensemble, ou bien qui, comme chez vous, restent petits ensemble ! » \end{center}

\begin{flushright}\textit{Le Gai Savoir}, Friedrich Nietzsche\end{flushright}

\begin{center}« \textitsl{Âne (ou Chameau).} — Ce sont les bêtes du désert (nihilisme). Ils portent, ils portent des fardeaux jusqu’au fond du désert. L’Âne a deux défauts : son Non est un faux non, un « non » du ressentiment. Et plus encore, son Oui (I-A, I-A) est un faux oui. Il croit qu’affirmer signifie porter, assumer. L’Âne est d’abord l’animal chrétien : il porte le poids des valeurs dites « supérieures à la vie ». Il affirme, mais n’affirme que les produits du nihilisme. »\end{center}

\begin{flushright}\textit{Nietzsche}, Gilles Deleuze\end{flushright}

\begin{center}« »\end{center}

\begin{center}« »\end{center}

\begin{center}« »\end{center}

\begin{center}« »\end{center}

\begin{center}« »\end{center}


\end{document}
